\documentclass[../../main.tex]{subfiles}
\begin{document}

{\bf [解答]}

% TODO: 概要図

題意より$O(0,0), A(1,1), B(2,2^n)$である．まず$f(x)$を求める．$O$，$A$を通るから$a \ (\neq 0)$を用いて，
\begin{align}
  f(x)-x=ax(x-1) \label{eq:1}
\end{align}
と書ける．$B$を通る条件から
\begin{align}
  2^n-2=2a \quad\therefore a=2^{n-1}-1 \label{eq:2}
\end{align}
であるから，\cref{eq:1}に代入して
\begin{align}
  f(x)=(2^{n-1}-1)x^2+(2-2^{n-1})x \label{eq:3}
\end{align}
である．二つの曲線の差分を$g_n(x)=x^n-f(x)$とおくと
\begin{align}
  g_n(x)
  &=x\left[x^{n-1}-(2^{n-1}-1)x-2+2^{n-1}\right] \nonumber\\
  &=x\left[x^{n-1}+x+2^{n-1}(1-x)-2\right] \nonumber\\
  &=x(x-1)h(x) \label{eq:4}
\end{align}
である．ただし
\begin{align}
  h(x)=(x^{n-2}+x^{n-3}+\cdots+x+2)-2^{n-1}
\end{align}
とおいた．$0\le x\le2$のとき，
\begin{align}
  h(x)
  &\le (2^{n-2}+2^{n-3}+\cdots+2+2)-2^{n-1} \\
  &=2^{n-1}-1+(1-2^{n-1})=0
\end{align}
であるから，\cref{eq:4}より$g_n(x)$の符号は
\begin{align}
  \begin{cases}
    g_n(x)\ge0&(0\le x\le1) \\
    g_n(x)\le0&(1\le x\le2)
  \end{cases}
\end{align}
となる．
従って題意の面積は
\begin{align}
  S_1=\int_0^1g_n(x)dx, \quad
  S_2=-\int_1^2g_n(x)dx
\end{align}
であり，
\begin{align}
  S_1=S_2\Longleftrightarrow \int_0^2g_n(x)dx=0 \\
  \Longleftrightarrow
  \left(\dfrac{1}{n+1}-\dfrac{1}{6}\right)2^{n+1}=\dfrac{4}{3}
\end{align}
をみたす$n$を求めればよい．$n\ge5$のときこの左辺は$0$以下になり不適だから，$n=3,4$が必要．実際に代入すると，適するのは$n=3$のときである．$\cdots$(答)
\newpage
\end{document}
